\chapter{Channels, Go-to-Market \& Sales}\label{ch:channels-gtm-and-sales}

% Scope: channel strategy (the 4th P), PLG vs sales-led GTM motion, sales-force sizing &
%   compensation, B2B selling & pipeline management. NEW chapter (closes the channels + sales/GTM gap).
% Sits upstream of the digital funnel (\cref{ch:digital-funnel-and-crm}): where you sell
%   determines what funnel you build.
% Cross-links: \cref{eq:cac}, \cref{eq:ltv-cac}, \cref{eq:viral-k}, \cref{eq:pipeline-ev},
%   \cref{eq:littles-law}, \cref{sec:kerr-trap}, \cref{eq:segment-score}.

\section{Channel Strategy \& the 4th P}\label{sec:channel-strategy}
% complexity: 4

How should a firm make its product physically and commercially available to customers?
Distribution --- the fourth ``P'' of the marketing mix --- determines who bears the cost of
reaching the customer, who owns the relationship, and ultimately how much of the unit margin
the firm retains. The decision is often treated as a logistics question, but it is primarily a
margin question: every intermediary added to the chain lowers the price the manufacturer
receives and, in return, offloads a portion of the reach, inventory, and service cost. The
strategic question is whether the margin given up is less than the fixed cost of reaching
those customers directly.

Let \(P\) be the end-user price, \(C_{\text{var}}\) the variable cost per unit, \(Q\) the
volume sold, and \(\Phi_{\text{fixed}}\) the direct-channel fixed overhead (a sales team,
a support organisation, transaction infrastructure). The contribution from each route is:
\begin{equation}\label{eq:channel-margin}
  \Pi_{\text{direct}}   = (P - C_{\text{var}})\,Q, \qquad
  \Pi_{\text{indirect}} = (P_w - C_{\text{var}})\,Q,
\end{equation}
where \(P_w < P\) is the wholesale price paid by the intermediary. After charging for the
direct fixed overhead, the firm prefers direct distribution when
\(\Pi_{\text{direct}} - \Phi_{\text{fixed}} > \Pi_{\text{indirect}}\), or equivalently when
\((P - P_w)\,Q > \Phi_{\text{fixed}}\): the unit-margin premium of going direct, scaled by
volume, must exceed the fixed cost of running the direct channel. At low volumes the
intermediary wins; at high volumes the direct route wins. The crossover is the channel
breakeven volume.

\begin{figure}[htbp]
  \centering
  \includegraphics[width=0.78\linewidth]{channel-structure}
  \caption{Three channel structures: direct (manufacturer to customer), one-level (via a
    single reseller or VAR), and two-level (via a distributor plus reseller). Each
    intermediary layer captures a margin spread from \cref{eq:channel-margin}, reducing the
    vendor's net unit revenue. The optimal structure minimises total distribution cost at
    the required market coverage.}
  \label{fig:channel-structure}
\end{figure}

\begin{example}[Direct vs.\ reseller channel decision]
The running SaaS company charges \money{50}/month (\money{600}/year). Its variable cost is
\money{29}/month (hosting \money{15} plus support, payment processing, and customer success
\money{14}), leaving a contribution margin of \money{21}/month on the direct channel.

A regional reseller proposes taking a \qty{30}{\percent} margin share on the end-user price,
so the vendor receives a wholesale price of \(P_w = \money{50}\times0.70 =
\money{35}/\)month. The reseller's indirect contribution margin per customer is then
\(\money{35} - \money{29} = \money{6}/\)month --- one-third of the direct margin. The margin
sacrifice is \(\money{21} - \money{6} = \money{15}\)/month per customer.

Suppose the firm must spend \money{90000}/month in direct sales and customer-success fixed
overhead to serve the reseller's territory without the reseller. The breakeven incremental
customer count at which direct becomes preferable is
\[
  Q^* \;=\; \frac{\Phi_{\text{fixed}}}{P - P_w} \;=\; \frac{\money{90000}}{\money{15}}
       \;=\; \num{6000}\ \text{customers}.
\]
The company currently serves roughly \num{6667} paying customers (\money{4000000} ARR
\(\div\) \money{600}/year). At that scale, the direct channel contributes
\(\num{6667}\times\money{21} = \money{140000}\)/month after variable costs, versus
\(\num{6667}\times\money{6} = \money{40000}\)/month indirect; the \money{90000} fixed-cost
gap is covered. The implication: the reseller arrangement made economic sense when the
installed base was below \num{6000}, but the firm has now crossed the threshold where
building a direct channel pays. If the reseller opens a new geography the firm cannot reach
at comparable cost, the indirect channel still creates value there --- the breakeven
calculation applies segment by segment, not company-wide.
\end{example}

\begin{admonition}{Warning}
Running both a direct sales team and a reseller in the same territory creates
\textit{dual-distribution channel conflict}: the reseller competes for the same prospect,
discounts to win, and ultimately trains customers to wait for the lower indirect price. The
firm then faces a choice between protecting the channel partner relationship and defending
the direct price. Channel agreements that define territory exclusivity or deal-registration
rights reduce this risk, but only if they are enforced consistently; inconsistent enforcement
is worse than no policy at all, because it creates uncertainty for every party.
\end{admonition}

\cref{eq:channel-margin} applies equally to physical goods and to software, but the
cost structure differs sharply: a SaaS firm's \(C_{\text{var}}\) is dominated by hosting
and customer-success labour, not inventory, so the fixed-cost threshold
\(\Phi_{\text{fixed}}\) --- the cost of a direct sales organisation --- dominates the
calculation at low volumes. That asymmetry is what makes the product-led motion examined in
\cref{sec:gtm-motion} so economically powerful: it builds a direct channel with near-zero
fixed cost per customer acquired. \textcite{coughlan2006marketing} develop the full
channel-design framework; \textcite{kotler2021} situate it within the integrated marketing mix.


\section{PLG vs.\ Sales-Led GTM Motion}\label{sec:gtm-motion}
% complexity: 4

Once the channel structure is settled, the firm must decide how demand enters the channel:
does the product sell itself through free trials and self-serve onboarding, or does a sales
team drive the entire journey from awareness to close? This is the go-to-market motion
decision, and it is one of the highest-leverage strategic choices an early-stage software
company makes, because it determines the shape of \cref{eq:cac}, the denominator of
\cref{eq:ltv-cac}, before the company has spent a dollar on acquisition.

In a \textit{product-led growth} (PLG) motion, the product itself is the primary acquisition
vehicle: users sign up, experience the core value within a self-serve trial, and convert to
paid without a human touch. CAC is low because there is no sales compensation embedded in
the acquisition cost; the viral coefficient (\cref{eq:viral-k}) can compound
the installed base at near-zero marginal cost when existing users invite colleagues. PLG
dominates when the product delivers immediate, tangible value to an individual user who can
make or influence a purchase decision autonomously --- typically SMB buyers with annual
contract values (ACV) below \money{5000}.

In a \textit{sales-led} motion, the product is gated: a prospect must engage a sales
representative before accessing a meaningful demonstration. CAC is high because it includes
base salary, commission, and the overhead of a sales development representative (SDR) who
qualifies the lead before handing it to an account executive. Sales-led is justified when
the decision involves multiple stakeholders, a security or legal review, or a procurement
process that a self-serve checkout cannot navigate --- typically enterprise ACV above
\money{50000}. \textcite{ross2011predictable}, the canonical practitioner reference for outbound SaaS sales, formalised the SDR / AE / CSM separation that this chapter takes as standard.

The hybrid motion, \textit{product-led sales} (PLS), layers a sales-assist team atop a
self-serve funnel. The productivity unlock is the \textit{product qualified lead} (PQL): an
active user who has crossed a demonstrable engagement threshold becomes the signal that
triggers SDR outreach. Because the sales rep contacts a user who has already experienced
value, the trust gap is closed before the first conversation. The segment-targeting logic of
\cref{eq:segment-score} maps directly onto the motion choice: segments with high LTV
potential but low self-serve conversion are PLS candidates; segments with high self-serve
conversion and low LTV are pure PLG.

\begin{figure}[htbp]
  \centering
  \includegraphics[width=0.86\linewidth]{plg-vs-slg}
  \caption{The PLG--SLG spectrum. Moving right, CAC rises, ACV rises, and the sales cycle
    lengthens. The PQL trigger (dashed line) marks the transition from self-serve to
    sales-assist; it is typically an in-product engagement milestone, not a marketing
    qualification score.}
  \label{fig:plg-vs-slg}
\end{figure}

\begin{example}[Hybrid PLG + SDR model]
The company currently acquires \num{200} customers per month through a self-serve funnel at
\(\text{CAC} = \money{600}\) (from \cref{eq:cac} in \cref{ch:customer-economics}).
Segmenting the user base reveals two distinct cohorts: SMB accounts that self-serve at
CAC\,\(\approx\)\,\money{300}, and mid-market accounts that stall before converting. An SDR
motion targeting mid-market PQLs carries CAC\,\(\approx\)\,\money{800}, but mid-market LTV
is roughly \(3\times\) that of SMB, so the \cref{eq:ltv-cac} ratio still clears the
3\(\times\) threshold.

The PQL trigger chosen is \textit{more than five reports generated in the first fourteen
days of the trial}. Users who hit this threshold convert at \qty{45}{\percent}; the baseline
conversion rate for all trial users is \qty{25}{\percent}. Routing a triggered user to an
SDR therefore captures an \qty{80}{\percent} relative lift in conversion probability before
the first human contact --- the product has already done the selling. The operational
implication: the product analytics pipeline must surface this signal in the CRM within
minutes of the threshold being crossed; a 24-hour lag loses roughly half the conversion
advantage as user attention moves elsewhere.
\end{example}

\begin{admonition}{Note}
The PLG versus sales-led decision is not a company-wide binary. Most companies above \money{2000000}
ARR operate both motions in parallel, segmented by buyer type. Forcing a single motion on
all segments misallocates resources: applying an SDR to a \money{600}/year SMB deal destroys
margin; leaving a \money{50000}/year enterprise prospect in a self-serve trial loses the deal to
a competitor with a human champion inside the account. The right diagnostic is the unit
economics of each motion by segment, not a philosophical stance on PLG.
\end{admonition}

The GTM motion feeds directly into the funnel mechanics of \cref{ch:digital-funnel-and-crm}:
a PLG motion fills the top of the funnel with trial activations measured by \cref{eq:viral-k};
a sales-led motion fills it with MQL handoffs from marketing; PLS uses both. The pipeline
valuation in \cref{eq:pipeline-ev} applies downstream regardless of motion, because
every deal in the CRM carries a face value and a stage probability whether it arrived via
self-serve or SDR outreach.


\section{Sales-Force Management \& Compensation}\label{sec:salesforce-management}
% complexity: 3

Building a sales team requires two separate models: a \textit{sizing model} that translates
a revenue target into a headcount requirement, and a \textit{compensation model} that aligns
the representative's incentives with the firm's value creation goals. Getting either wrong
is expensive. An undersized team leaves ARR on the table; an oversized team inflates
sales-and-marketing expense without proportional revenue. A poorly designed commission plan
rewards the wrong behaviour and, through Kerr's folly (\cref{sec:kerr-trap} in
\cref{ch:organizational-behavior}), systematically produces outcomes the
firm did not intend.

The headcount formula flows directly from the ARR target, the quota per account executive,
and the historical attainment rate:
\begin{equation}\label{eq:sales-productivity}
  \text{reps} \;=\;
    \left\lceil\frac{\text{ARR target}}{\text{quota}\times\text{attainment}}\right\rceil.
\end{equation}
The ceiling function ensures the firm fields enough capacity to reach the target even when
reps fall short of \qty{100}{\percent} attainment. The attainment factor is the critical
input: a plan built on \qty{100}{\percent} attainment will miss every quarter; the
historical blended rate, typically \qty{70}{\percent}--\qty{85}{\percent} for a mature SaaS
team, is the right denominator.

Compensation is structured as an on-target earnings (OTE) split between a fixed base and
performance-linked variable:
\begin{equation}\label{eq:ote-structure}
  \text{OTE} \;=\; \text{base} \;+\; \text{variable}, \qquad
  \text{variable} \;=\; \text{quota} \times \text{commission rate}.
\end{equation}
A \qty{50}{\percent}/\qty{50}{\percent} base-to-variable split is standard for an
enterprise AE role; inside sales and SDR roles run closer to \qty{70}{\percent}/\qty{30}{\percent}
because the individual's influence over deal size is smaller. The commission rate in
\cref{eq:ote-structure} is simply the variable target divided by the quota, and it sets the
slope of the incentive curve.

\begin{example}[Sizing and costing the sales team for \money{4000000}\,$\to$\,\money{6000000} ARR]
The firm targets \money{6000000} ARR against a current base of \money{4000000}, requiring
\money{2000000} in new ARR. Each account executive carries a \money{500000} quota; at the
historical \qty{80}{\percent} blended attainment rate, \cref{eq:sales-productivity} gives:
\[
  \text{reps} \;=\;
    \left\lceil\frac{\money{2000000}}{\money{500000}\times0.80}\right\rceil
    \;=\; \lceil 5.00 \rceil \;=\; 5\ \text{AEs}.
\]
Each AE has OTE of \money{120000} on a \qty{50}{\percent}/\qty{50}{\percent} split: base
\money{60000}, variable \money{60000}. From \cref{eq:ote-structure} the commission rate is
\(\money{60000}/\money{500000} = \qty{12}{\percent}\) of ARR booked. Five AEs cost
\money{600000} in total compensation --- the incremental sales-and-marketing line for this
growth tranche.

Does the investment clear the LTV:CAC hurdle? At \qty{80}{\percent} attainment, five AEs
produce \(5\times\money{500000}\times0.80=\money{2000000}\) in new ARR, or roughly
\num{3333} new customers at the \money{600}/year ACV from the fact sheet. The blended
sales-assisted CAC is \(\money{600000}/\num{3333}\approx\money{180}\). Against the CLV of
\money{3150} from \cref{ch:customer-economics}, the LTV:CAC ratio is
\(\money{3150}/\money{180}\approx17.5\times\), comfortably above the 3\(\times\) floor of
\cref{eq:ltv-cac}. The binding constraint is therefore not economics but execution:
finding five AEs who hit quota.
\end{example}

\begin{admonition}{Warning}
Accelerator clauses that multiply commissions above \qty{100}{\percent} quota attainment
create a rational incentive for representatives to sandbag early in the quarter --- hoarding
deals until they have enough pipeline to guarantee hitting the accelerator threshold before
submitting them. The firm sees lumpy, end-of-quarter booking patterns that compress the
implementation and onboarding queue, inflate churn in the following quarter, and make the
forecast unreliable. The structural fix is to pay commission on \textit{retained} ARR at the
end of the first renewal period rather than on booked ARR at signature: a representative
whose customer churns within twelve months clawbacks a portion of the initial commission.
This directly addresses the Kerr Trap (\cref{sec:kerr-trap}): paying on booked ARR rewards
closing, while hoping for renewal retention. \textcite{dixon2011challenger} document
the full taxonomy of sales-behaviour distortions that emerge from misaligned commission design.
\end{admonition}

The sales productivity formula and the OTE structure are inputs to the operating model in
\cref{ch:operating-model}, where they appear as line items in the sales-and-marketing
expense projection. Connecting compensation design to churn ensures the ARR growth modelled
in \cref{eq:mrr-growth} is not eroded by sales-induced early churn.


\section{B2B Selling \& Pipeline Management}\label{sec:b2b-selling}
% complexity: 3

A sales team that cannot forecast its own revenue is a liability, not an asset. Raw pipeline
--- the sum of all open deal face values --- is the canonical vanity metric of enterprise
sales: it treats a deal in the earliest discovery call as equivalent to a deal in final legal
review. Weighted pipeline replaces the face value with the expected value of each
opportunity, where the expectation is conditioned on where the deal sits in the stage
funnel:
\begin{equation}\label{eq:pipeline-forecast}
  \widehat{\text{ARR}} \;=\; \sum_{i=1}^{n} v_i \cdot p_{\,\text{stage}(i)},
\end{equation}
where \(v_i\) is the face value of deal \(i\) and \(p_{\,\text{stage}(i)}\) is the
historical close probability for its current stage. \Cref{eq:pipeline-forecast} is the
ARR analogue of \cref{eq:pipeline-ev} in \cref{ch:digital-funnel-and-crm}: both apply the
expected-value operator to a portfolio of probabilistic conversions. The stage probabilities
must be calibrated to observed win rates at each stage, not assigned by managerial optimism;
a pipeline model built on inflated stage probabilities fails at the same rate as no model at
all, with the additional harm of creating false confidence.

Alongside the expected-value snapshot, Little's Law from \cref{eq:littles-law} gives the
average sales-cycle length from pipeline inventory and throughput:
\[
  W \;=\; \frac{L}{\lambda},
\]
where \(L\) is the number of open deals, \(\lambda\) is the number of deals closed per
period, and \(W\) is the average time a deal spends in the pipeline. A rising
\(W\) is an early-warning signal that deals are stalling, not that pipeline is healthy.

The SPIN questioning sequence (\textcite{rackham1988spin}) and the Challenger methodology
(\textcite{dixon2011challenger}) operate at the level of the individual sales call but
produce the stage progression that fills \cref{eq:pipeline-forecast}. SPIN's
situation--problem--implication--need-payoff arc surfaces the economic pain that justifies
a purchase decision; the Challenger's teach--tailor--take-control arc reframes the buyer's
mental model before a competitor can anchor it. Both frameworks exist to move deals
rightward through the pipeline stages, increasing \(p_{\,\text{stage}(i)}\) for the
in-flight portfolio.

\begin{example}[Weighted pipeline and Little's Law]
The sales team carries \num{40} open deals distributed evenly across four pipeline stages:
\num{10} in Discovery, \num{10} in Demo, \num{10} in Proposal, and \num{10} in Negotiation.
Each deal has a face value of \money{8250} ARR. Historical stage close-rates are
\qty{15}{\percent} at Discovery, \qty{35}{\percent} at Demo, \qty{60}{\percent} at Proposal,
and \qty{85}{\percent} at Negotiation. Applying \cref{eq:pipeline-forecast}:
\begin{align*}
  \widehat{\text{ARR}}
    &= 10\times\money{8250}\times0.15
    \;+\; 10\times\money{8250}\times0.35 \\
    &\quad +\; 10\times\money{8250}\times0.60
    \;+\; 10\times\money{8250}\times0.85 \\
    &= \money{12375} + \money{28875} + \money{49500} + \money{70125}
     \;\approx\; \money{161000}.
\end{align*}
The raw pipeline face value is \(40\times\money{8250}=\money{330000}\) --- more than double
the \money{161000} expected ARR. A sales manager reporting the raw number as ``pipeline
coverage'' would overstate the true forecast by \qty{105}{\percent}.

Applying Little's Law: if the team closes \num{10} deals per month, the average deal spends
\(W = 40/10 = 4\) months in the pipeline. If the team adds \num{5} SDR-sourced opportunities
per month but closes at the same \num{10}/month rate, \(L\) grows to \num{60} within four
months and \(W\) extends to \(60/10 = 6\) months. The pipeline looks healthier (more deals),
but the expected close date for any given deal has slipped by two months --- a forecast
problem masquerading as a pipeline success.
\end{example}

\begin{admonition}{Note}
Qualification frameworks --- such as MEDDPICC (Metrics, Economic Buyer, Decision Criteria,
Decision Process, Paper Process, Identify Pain, Champion, Competition) --- serve the same
function for the pipeline that segment-scoring (\cref{eq:segment-score}) serves for the
addressable market: they force an explicit, measurable estimate of opportunity quality before
committing sales capacity. An opportunity that fails the qualification checklist should be
de-staged or disqualified rather than carried at an optimistic probability; the cost is a
temporarily smaller pipeline, and the benefit is a \cref{eq:pipeline-forecast} that
management can actually rely on.
\end{admonition}

Channel structure, GTM motion, sales-force design, and pipeline management are the four
decision layers that convert a product and a target segment into predictable revenue.
Each layer feeds the next: the channel determines the acquisition cost structure
(\cref{eq:channel-margin}), the motion determines whether that cost is dominated by fixed
sales overhead or variable product investment, the compensation design determines whether
the reps' behaviour aligns with LTV maximisation (\cref{eq:ltv-cac}), and the pipeline
model determines whether the organisation can read its own progress. All four feed the
digital funnel in \cref{ch:digital-funnel-and-crm} and ultimately the revenue line of the
operating model in \cref{ch:operating-model}.
